\subsubsection{Random Forest}

Como hemos comentado sobre los Random Forest, estos construyen un conjunto o ''bosque'' de árboles de decisión donde cada árbol se entrena de forma ligeramente diferente, introduciendo aleatoriedad controlada.

\subsubsection*{Justificación del modelo}

Se seleccionó Random Forest como uno de los clasificadores base debido a su robustez frente al sobreajuste, su capacidad para manejar características heterogéneas y su interpretabilidad mediante el análisis de importancia de variables.

\subsubsection*{Configuración experimental}

El proceso experimental para Random Forest se estructuró en tres fases: (1) establecimiento de un modelo baseline, (2) validación cruzada para evaluar robustez, (3) optimización de hiperparámetros mediante búsqueda sistemática.

\paragraph{División de datos}
El conjunto de datos se particionó en entrenamiento (80\%) y prueba (20\%) utilizando muestreo estratificado para mantener la proporción de clases en ambos conjuntos. Esta estratificación es crucial cuando existe desbalance entre voces reales y generadas.

\begin{lstlisting}[language=Python, caption=División estratificada de datos]
X_train, X_test, y_train, y_test = train_test_split(
    X, y, 
    test_size=0.2, 
    random_state=12, 
    stratify=y  # Mantiene proporcion de clases
)
\end{lstlisting}

\paragraph{Modelo baseline}
Se estableció una configuración inicial con hiperparámetros por defecto para obtener una referencia de rendimiento:

\begin{lstlisting}[language=Python, caption=Configuración del modelo baseline Random Forest]
rf_baseline = RandomForestClassifier(
    n_estimators=100,      # Numero de arboles
    max_depth=10,           # Profundidad maxima
    min_samples_split=5,    # Minimo muestras para dividir nodo
    min_samples_leaf=2,     # Minimo muestras en hoja
    max_features='sqrt',    # Caracteristicas por division
    random_state=12,
    n_jobs=-1,              # Usar todos los cores
    class_weight='balanced' # Ajuste para clases desbalanceadas
)
\end{lstlisting}

El parámetro \texttt{class\_weight='balanced'} ajusta los pesos de las clases de forma inversamente proporcional a su frecuencia, compensando posibles desbalances en el conjunto de datos.

\subsubsection*{Validación cruzada}

Para evaluar la robustez del modelo y evitar que los resultados dependan de una única partición de los datos, se realizó validación cruzada estratificada con 5 folds:

 \begin{lstlisting}[language=Python, caption=Validación cruzada con Random Forest]
from sklearn.model_selection import StratifiedKFold, cross_val_score

cv = StratifiedKFold(n_splits=5, shuffle=True, random_state=12)

cv_scores_accuracy = cross_val_score(rf_baseline, X_train, y_train, 
                                     cv=cv, scoring='accuracy')
cv_scores_f1 = cross_val_score(rf_baseline, X_train, y_train, 
                               cv=cv, scoring='f1')
cv_scores_auc = cross_val_score(rf_baseline, X_train, y_train, 
                                cv=cv, scoring='roc_auc')
\end{lstlisting}

La validación cruzada proporciona una estimación más realista del rendimiento esperado en datos no vistos, reportándose la media y la desviación estándar de las métricas.

\subsubsection*{Optimización de hiperparámetros}

Se realizó una búsqueda sistemática de hiperparámetros mediante \texttt{GridSearchCV} con validación cruzada de 5 folds, optimizando el F1-score. El espacio de búsqueda incluyó:

\begin{lstlisting}[language=Python, caption=Grid Search para Random Forest]
param_grid = {
    'n_estimators': [50, 100, 200],	# Numero de arboles en el bosque
    'max_depth': [5, 10, 15, None],	# Profundidad maxima de cada arbol
    'min_samples_split': [2, 5, 10],	# Minimo de muestras para dividir un nodo
    'min_samples_leaf': [1, 2, 4],	# Minimo de muestras en una hoja
    'max_features': ['sqrt', 'log2']	# Numero de caracteristicas a considerar por division
}

grid_search = GridSearchCV(
    RandomForestClassifier(random_state=12, class_weight='balanced'),
    param_grid,
    cv=cv,
    scoring='f1',  # Optimizacion de F1-score
    n_jobs=-1,
    verbose=1
)

grid_search.fit(X_train, y_train)
rf_optimized = grid_search.best_estimator_
\end{lstlisting}

\subsubsection*{Métricas de evaluación}

Para una evaluación exhaustiva de los modelos, se emplearon las siguientes métricas:

\begin{itemize}
    \item \textbf{Accuracy (Exactitud)}: Proporción de predicciones correctas sobre el total.
    \[
    \text{Accuracy} = \frac{TP + TN}{TP + TN + FP + FN}
    \]
    
    \item \textbf{Precision (Precisión)}: Proporción de predicciones positivas correctas.
    \[
    \text{Precision} = \frac{TP}{TP + FP}
    \]
    
    \item \textbf{Recall (Sensibilidad)}: Proporción de positivos reales correctamente identificados.
    \[
    \text{Recall} = \frac{TP}{TP + FN}
    \]
    
    \item \textbf{F1-Score}: Media armónica de precisión y recall.
    \[
    F1 = 2 \cdot \frac{\text{Precision} \cdot \text{Recall}}{\text{Precision} + \text{Recall}}
    \]
    
    \item \textbf{AUC-ROC} (Área bajo la curva ROC): Mide la capacidad del modelo para distinguir entre clases, independientemente del umbral de clasificación. Un valor de 0.5 indica clasificación aleatoria, mientras que 1.0 indica clasificación perfecta.
    
    \item \textbf{MCC} (Coeficiente de correlación de Matthews): Métrica equilibrada que considera todos los elementos de la matriz de confusión, especialmente útil en problemas con clases desbalanceadas. Su rango es \([-1, 1]\), donde 1 indica predicción perfecta, 0 predicción aleatoria y -1 predicción inversa perfecta.
    \[
    \text{MCC} = \frac{TP \cdot TN - FP \cdot FN}{\sqrt{(TP+FP)(TP+FN)(TN+FP)(TN+FN)}}
    \]
\end{itemize}

La función de evaluación implementada automatiza el cálculo de todas estas métricas y genera visualizaciones:

\begin{lstlisting}[language=Python, caption=Evaluación de un modelo]
def evaluacion_modelo(model_name, y_true, y_pred, y_pred_proba):
    """
    Evalua un modelo con metricas completas y visualizaciones.
    """
    # Calculo de metricas
    accuracy = accuracy_score(y_true, y_pred)
    auc = roc_auc_score(y_true, y_pred_proba)
    f1 = f1_score(y_true, y_pred)
    mcc = matthews_corrcoef(y_true, y_pred)
    precision = precision_score(y_true, y_pred)
    recall = recall_score(y_true, y_pred)
    
    # Matriz de confusion
    cm = confusion_matrix(y_true, y_pred)
    
    # Visualizaciones (matriz de confusion, curva ROC)
    # ...
    
    return {
        'accuracy': accuracy, 'precision': precision,
        'recall': recall, 'f1': f1,
        'auc': auc, 'mcc': mcc,
        'confusion_matrix': cm
    }
\end{lstlisting}